% ============================================================================
%  Reusable preamble for tight, black-&-white LaTeX teaching-notes PDFs.
%  Copy this into your main.tex (or \input it). Set \figpath to your figures dir.
%  Design goals: regular tex-to-pdf look, B&W, information-dense, page numbers,
%  bounded (unbreakable) KEY POINT / NOTE / RESULT / DEFINITION boxes.
% ============================================================================
\documentclass[11pt]{article}

% ---------- fonts & encoding ----------
\usepackage[T1]{fontenc}
\usepackage{lmodern}
\usepackage[utf8]{inputenc}
\usepackage{textcomp}
\usepackage{amsmath,amssymb}
\usepackage{microtype}

% ---------- tight, information-dense page ----------
\usepackage[a4paper,margin=2.0cm,headsep=8pt,footskip=20pt]{geometry}
\setlength{\parindent}{0pt}
\setlength{\parskip}{4pt plus 1pt minus 1pt}
\linespread{1.02}

% ---------- graphics, tables, lists ----------
\usepackage{graphicx}
\usepackage{booktabs}
\usepackage{array}
\usepackage{tabularx}
\usepackage{longtable}
\usepackage{ragged2e}
\usepackage[inline]{enumitem}
\setlist{nosep,leftmargin=1.5em,topsep=2pt,itemsep=2pt}
\usepackage{caption}
\captionsetup{font=small,labelfont=bf,skip=4pt}
\usepackage{float}          % provides [H] so large figures/tables stay in-flow

% tighter tables everywhere
\setlength{\tabcolsep}{4pt}
\renewcommand{\arraystretch}{1.06}

% ---------- black & white palette (colour names resolve to black) ----------
\usepackage{xcolor}
\definecolor{navy}{gray}{0}\definecolor{teal}{gray}{0}
\definecolor{dgreen}{gray}{0}\definecolor{dred}{gray}{0}\definecolor{ochre}{gray}{0}
\usepackage[most]{tcolorbox}
\usepackage{tikz}
\usetikzlibrary{arrows.meta,calc}

% ---------- compact, black section headings ----------
\usepackage{titlesec}
\titleformat{\section}{\normalfont\large\bfseries}{\thesection}{0.6em}{}
\titleformat{\subsection}{\normalfont\normalsize\bfseries}{\thesubsection}{0.6em}{}
\titlespacing*{\section}{0pt}{9pt plus 1pt minus 1pt}{3pt}
\titlespacing*{\subsection}{0pt}{6pt plus 1pt minus 1pt}{2pt}

% ---------- header / footer with page numbers on every page ----------
\usepackage{fancyhdr}
\pagestyle{fancy}\fancyhf{}
\renewcommand{\headrulewidth}{0.4pt}
\fancyhead[L]{\small\itshape \runningtitle}   % define \runningtitle before \begin{document}
\fancyhead[R]{\small\leftmark}
\fancyfoot[C]{\small\thepage}
\renewcommand{\sectionmark}[1]{\markboth{\thesection.\ #1}{}}

\usepackage[hidelinks]{hyperref}
\usepackage{cleveref}

% ============ teaching boxes: black & white, bounded, UNBREAKABLE ============
% Unbreakable so a box never splits across a page (title + one orphan line).
% If it doesn't fit, it moves whole to the next page.
\tcbset{bwbox/.style={enhanced,colback=white,colframe=black,
  colbacktitle=black!10,coltitle=black,boxrule=0.6pt,titlerule=0pt,
  left=6pt,right=6pt,top=3pt,bottom=3pt,toptitle=1.5pt,bottomtitle=1.5pt,
  arc=1pt,fonttitle=\bfseries\footnotesize,before skip=6pt,after skip=6pt}}
\newtcolorbox{keynote}[1][]{bwbox,title={KEY POINT},#1}      % important takeaway
\newtcolorbox{notebox}[1][]{bwbox,title={NOTE},#1}           % side aid / derivation note
\newtcolorbox{resultbox}[1][]{bwbox,boxrule=1.1pt,title={RESULT \& MEANING},#1} % final result
\newtcolorbox{defbox}[1][]{bwbox,title={DEFINITION},#1}      % boxed definition
\newtcolorbox{oneline}{enhanced,colback=black!5,colframe=black,boxrule=0.6pt,
  left=6pt,right=6pt,top=3pt,bottom=3pt,arc=1pt}            % one-line "in a nutshell"

% set these per document, before \begin{document}:
% \newcommand{\figpath}{/abs/path/to/figs}
% \newcommand{\runningtitle}{Short Running Title}
